\chapter{DISCUSSION AND CONCLUSION}
\onehalfspacing
\label{chap:discussion}
\vspace{-3em}
\section{Thesis Outcome}
This study yielded four key outcomes: first, development of an LLM-based screening tool for borderline personality disorder; second, creation of an expert-annotated QnA dataset tailored to BPD screening; third, execution of a comprehensive comparative analysis across machine learning, deep learning, and large language models; and fourth, identification and deployment of ELECTRA trained on augmented text—as the optimal model, achieving an 86.84 \% F1-score in the final screening tool.

\section{Thesis Implication and Contribution}
 The LLM-based screening tool will be a powerful addition to psychologists' and psychiatrists' methods in screening their patients for borderline personality disorder screening. The general population can also use this tool to self-assess existence of borderline personality disorder symptoms, and seek appointments with specialists after using the tool for a definitive diagnosis.

The contributions of this research are threefold. First, a novel QnA classification dataset is formulated for detecting borderline personality disorder symptoms. Second, an LLM-based screening tool is developed to assist psychologists and psychiatrists in making more informed diagnoses. Finally, a new method is proposed for explainable and detailed screening of borderline personality disorder, improving upon traditional dichotomous questions and Likert scales.

\section{Thesis Limitation}

It is crucial to acknowledge the limitations of this thesis. The following are some of the shortcomings of this thesis:
\begin{itemize}
  \item This tool is specifically tailored to detect borderline personality disorder, and is not applicable to other personality disorders
  \item Leveraging sequence-to-sequence LLMS in classification of mental health QnA pairs is yet to be investigated
  \item A better approach for determining criteria B-F for general personality disorder needs to be studied
\end{itemize}

\section{Future Work}

The future work of this research is to explore more ensembling methods - bagging, boosting, blending - to obtain potential increase in accuracies. Moreover, a method to incorporate future responses into further improving the deployed LLM. Possible inclusion of speech patterns when collecting responses is another avenue to be explored, which may enhance screening accuracy. Furthermore, investigating more ways to solve class imbalances, such as weighted focal loss, combination of undersampling and oversampling can be carried out in future. Modifying the BERT architecture to better handle question-answer pairs is yet another avenue to explore in future.  


\endinput



 

